\begin{ouverture}[calcul]{5}{Expressions, distributivité et preuves}{Partie 2 · Algèbre et modèles}
{Formules · Développer · Factoriser · Démontrer}
\objectif{Produire une formule et lire la structure d’une expression.}
\objectif{Développer et factoriser par distributivité simple.}
\objectif{Utiliser le calcul littéral pour démontrer une propriété générale.}
\prerequis{Produits de relatifs et priorités opératoires.}
\end{ouverture}
\pagestyle{otMain}\essentiel\fontsize{11}{15}\selectfont
\phantomsection\label{cours-5-1}
\essSection{1}{Une lettre pour plusieurs valeurs}
\begin{definition}Une expression littérale contient des lettres qui représentent des nombres. $3x$ signifie $3\times x$, et $x^2$ signifie $x\times x$. Calculer la valeur d’une expression consiste à remplacer chaque lettre par la valeur donnée, en conservant des parenthèses autour d’un nombre négatif.\end{definition}
\begin{exemple}Pour $x=-2$, $3x^2-5x=3\times(-2)^2-5\times(-2)=12+10=22$. $3x+4$ est une somme ; $3(x+4)$ est un produit. L’opération effectuée en dernier détermine la structure.\end{exemple}
\begin{propriete}Réduire consiste à regrouper les termes de même nature : $5x+2x-3=7x-3$, $x\times x=x^2$ et $3x\times2x=6x^2$. En revanche, $x$ et $x^2$ ne sont pas des termes de même nature : $x+x^2$ ne se réduit pas.\end{propriete}
\begin{definition}[Écrire un nombre quelconque]Pour tout entier $n$, $2n$ est un nombre pair et $2n+1$ un nombre impair ; $n$ et $n+1$ sont deux entiers consécutifs. Une formule produite doit être testée sur un exemple pour contrôler sa vraisemblance.\end{definition}
\begin{attention}Tester une égalité pour quelques nombres peut trouver une erreur, mais ne prouve pas qu’elle est vraie pour toutes les valeurs. Pour une preuve générale, on transforme les expressions avec des propriétés vraies pour tous les nombres.\end{attention}
\clearpage
\phantomsection\label{cours-5-2}
\essSection{2}{La distributivité dans les deux sens}
\begin{propriete}Pour tous nombres $k$, $a$, $b$ : $k(a+b)=ka+kb$ et $k(a-b)=ka-kb$, ainsi que $(a+b)k=ak+bk$. Passer du produit à la somme s’appelle développer. Passer de la somme au produit s’appelle factoriser ; on repère alors un facteur commun.\end{propriete}
\begin{exemple}$-3(2x-5)=-6x+15$. Le facteur $-1$ explique l’opposé d’une expression : $-(4x-2)=-1\times(4x-2)=-4x+2$. Dans l’autre sens, $12x^2+8x=4x(3x+2)$ et $x^2-x=x(x-1)$.\end{exemple}
\label{fig-cours-5-2-cours-distributivite}\figureGeo{cours-distributivite}

\begin{savoirfaire}[Démontrer avec une lettre]Choisir une lettre pour le nombre de départ. Écrire l’expression de chaque programme, puis développer et réduire. Obtenir la même expression prouve qu’ils donnent le même résultat pour tout nombre de départ.\end{savoirfaire}
\begin{attention}On simplifie des facteurs d’un produit, pas des termes séparés par une addition. Les produits de deux parenthèses et les identités remarquables seront étudiés en Troisième.\end{attention}

\clearpage\rubrique{Découverte}
\begin{blocLivre}{ex-5-1}\begin{exercice}[1]{Lire la structure}
Dire si $5x+15$, $5(x+3)$ et $x(2x-1)$ sont des sommes ou des produits.
\end{exercice}\end{blocLivre}
\begin{blocLivre}{ex-5-2}\begin{exercice}[1]{Substituer un négatif}
Calculer $A=4x-7$ pour $x=-3$.
\end{exercice}\end{blocLivre}
\begin{blocLivre}{ex-5-3}\begin{exercice}[1]{Réduire}
Réduire $8x-3+2x+7$, puis $3x\times2x+x^2$.
\end{exercice}\end{blocLivre}
\begin{blocLivre}{ex-5-4}\begin{exercice}[1]{Supprimer des parenthèses}
Écrire sans parenthèses puis réduire : $-(5-4x)$ et $2x-(3x-1)$.
\end{exercice}\end{blocLivre}
\begin{blocLivre}{ex-5-5}\begin{exercice}[1]{Factoriser}
Factoriser $18x+30$ avec le plus grand facteur entier commun.
\end{exercice}\end{blocLivre}

\rubrique{Entraînement}
\begin{blocLivre}{ex-5-6}\begin{exercice}[2]{Distribuer un signe négatif}
Développer et réduire $-2(5x-4)+3x$.
\end{exercice}\end{blocLivre}
\begin{blocLivre}{ex-5-7}\begin{exercice}[2]{Un facteur littéral}
Développer et réduire $3x(2x+5)-4x$.
\end{exercice}\end{blocLivre}
\begin{blocLivre}{ex-5-8}\begin{exercice}[2]{Factoriser avec une lettre}
Factoriser $15x^2-10x$.
\end{exercice}\end{blocLivre}
\begin{blocLivre}{ex-5-9}\begin{exercice}[2]{Équivalence de programmes}
Programme A : ajouter 4 à un nombre, puis multiplier par 3. Programme B : tripler le nombre, puis ajouter 12. Sont-ils équivalents ?
\end{exercice}\end{blocLivre}
\begin{blocLivre}{ex-5-10}\begin{exercice}[2]{Réparer une réduction}
Un élève écrit $4x+3x^2=7x^3$. Expliquer son erreur et donner la valeur de l’expression pour $x=2$.
\end{exercice}\end{blocLivre}

\rubrique{Problèmes}
\begin{blocLivre}{ex-5-11}\begin{exercice}[3]{Une aire composée}
Un rectangle a pour hauteur 4 cm et pour largeur $x+3$ cm, avec $x>0$. Exprimer son aire de deux façons.
\label{fig-ex-5-11}\figureGeo{aire-distrib}
\end{exercice}\end{blocLivre}
\begin{blocLivre}{ex-5-12}\begin{exercice}[3]{Un périmètre variable}
Un rectangle mesure $2x+1$ cm sur $x+3$ cm, avec $x>0$. Exprimer puis réduire son périmètre. Le calculer pour $x=4$.
\label{fig-ex-5-12}\figureGeo{rectangle-litteral}
\end{exercice}\end{blocLivre}
\begin{blocLivre}{ex-5-13}\begin{exercice}[3]{Un calcul mental justifié}
Calculer $39\times27+61\times27$ sans poser deux multiplications. Expliquer la propriété utilisée.
\end{exercice}\end{blocLivre}
\begin{blocLivre}{ex-5-14}\begin{exercice}[3]{Trois entiers consécutifs}
Montrer que la somme de trois entiers consécutifs est toujours un multiple de 3. Que vaut-elle quand le premier entier est 41 ?
\end{exercice}\end{blocLivre}
\begin{blocLivre}{ex-5-15}\begin{exercice}[3]{Une différence constante}
Programme A : multiplier le nombre de départ par 5 puis retirer 7. Programme B : retirer 2 au nombre, multiplier par 5 puis ajouter 6. Montrer que l’écart A moins B est constant.
\end{exercice}\end{blocLivre}

\clearpage
